\begin{ouverture}[college3]{12}{L’énergie : transferts et conversions}{Partie 3 · L’énergie, ses transferts et ses conversions}{Où va l’énergie d’une lampe ?}
\objectif{Distinguer source, forme, transfert et conversion.}
\objectif{Construire une chaîne et un bilan énergétiques simples.}
\objectif{Relier rayonnement, climat et choix de ressources.}
\prerequis{Combustions, mouvement et lecture d’un schéma.}
\end{ouverture}
\pcPage{Le cours}{L’énergie : transferts et conversions}
\phantomsection\label{ess-12}
\essSection{1}{Nommer avec précision}
Une source fournit de l’énergie à un système. Une forme décrit comment elle est stockée ou se manifeste : énergie chimique d’une pile, cinétique d’un objet en mouvement, potentielle liée à une position, thermique ou nucléaire. Le Soleil fournit un rayonnement ; une éolienne reçoit de l’énergie du vent. Le mot « électrique » qualifie notamment un transfert électrique dans une chaîne. Un dispositif convertit l’énergie sans la créer.\par
\pcFigure{figures/c12-chaine.pdf}{Pour la lampe : transfert électrique entrant, transferts par rayonnement et thermiques sortants. Les flèches ne représentent pas de la matière.}
\essSection{2}{Transférer et convertir}
Un transfert fait passer de l’énergie d’un système à un autre : électriquement, mécaniquement, thermiquement ou par rayonnement. Une conversion change la forme sous laquelle on décrit l’énergie dans le fonctionnement étudié. Un moteur électrique reçoit de l’énergie électriquement et met une charge en mouvement, tout en échauffant son environnement. Une chaîne nomme les systèmes et le type de chaque transfert.\par
\clearpage
\essSection{3}{Un bilan qui tient compte de tout}
L’énergie se conserve. Pour un système, la variation de l’énergie stockée correspond à l’énergie reçue moins l’énergie cédée. Si un convertisseur fonctionne sans variation de son stockage, l’énergie entrante égale la somme des énergies sortantes. Les « pertes » sont des transferts qui ne répondent pas à l’effet utile recherché ; elles ne constituent pas une disparition. L’énergie s’exprime en joules (J).\par
\pcFigure{figures/c12-terre.pdf}{Bilan qualitatif Terre-atmosphère : rayonnement solaire reçu, part renvoyée et rayonnement thermique émis. Les flèches ne sont pas quantitatives.}
\essSection{4}{Ressources et climat}
Le charbon, le pétrole et le gaz sont des ressources fossiles ; leur combustion émet du dioxyde de carbone. Une ressource renouvelable se renouvelle à une échelle compatible avec l’usage humain, sous conditions d’exploitation. « Renouvelable » ne signifie pas « sans impact ». Certains gaz à effet de serre absorbent et émettent du rayonnement infrarouge terrestre. Si la Terre reçoit plus d’énergie qu’elle n’en cède, son énergie stockée augmente, ce qui peut élever sa température moyenne.\par
\begin{pcEncadre}{method}{Un exemple expliqué}Une lampe reçoit $100\unite J$ électriquement et cède $25\unite J$ par rayonnement utile dans un modèle sans stockage. Les autres transferts totalisent $100-25=75\unite J$. L’énergie totale se conserve, mais toute l’énergie reçue n’est pas transférée sous la forme utile à l’éclairage.\end{pcEncadre}
\pcPage{Expérimenter et retenir}{L’énergie : transferts et conversions}
\phantomsection\label{activite-12}
\begin{pcEncadre}{method}{Investigation en classe}\textbf{Mission.} Étudier une lampe scolaire, un petit moteur et une cellule photovoltaïque, sur documents ou matériel autorisé. Définir chaque système, repérer l’entrée et les sorties puis construire une chaîne. Chercher un indice d’un transfert moins visible, par exemple un échauffement mesuré avec une sonde. Ne pas toucher une pièce chaude. \textbf{À rendre :} chaîne légendée, effet utile choisi, limites et bilan chiffré seulement lorsque toutes les données nécessaires sont fournies.\end{pcEncadre}
\pcSousTitre{Ce que je dois savoir expliquer}\begin{questions}
\item Une source n’est pas la même chose qu’une forme d’énergie.
\item Un transfert relie deux systèmes ; une conversion décrit leur fonctionnement.
\item L’énergie totale se conserve lorsqu’on inclut tous les transferts et stockages.
\item Un transfert non utile ne fait pas disparaître l’énergie.
\item Le rayonnement permet un transfert à travers le vide.
\end{questions}
\begin{pcEncadre}{attention}{Une erreur à éviter}« Le moteur fabrique de l’énergie. » Il convertit de l’énergie reçue ; la production usuelle d’électricité désigne une conversion et un transfert.\end{pcEncadre}
\pcRepere{cor-activite-12}{Pistes pour l’investigation}\pcGrille
\clearpage\markright{Je démarre}
\rubrique{Je démarre}
\par\noindent\begin{minipage}{\linewidth}
\begin{exercice}[1]{Identifier une source}
\begin{questions}
\item Dans une lampe alimentée par une pile, quel système fournit l’énergie à la lampe ?
\item Quelle énergie est stockée dans la pile avant utilisation ?
\end{questions}
\end{exercice}
\end{minipage}
\par\noindent\begin{minipage}{\linewidth}
\begin{exercice}[1]{Nommer une forme}
\begin{questions}
\item Quelle forme associe-t-on au mouvement d’un vélo ?
\item Et à un objet placé en hauteur dans le champ de pesanteur ?
\end{questions}
\end{exercice}
\end{minipage}
\par\noindent\begin{minipage}{\linewidth}
\begin{exercice}[1]{Une unité}
\begin{questions}
\item Quelle est l’unité SI de l’énergie ?
\item Le watt est-il une autre écriture du joule ?
\end{questions}
\end{exercice}
\end{minipage}
\par\noindent\begin{minipage}{\linewidth}
\begin{exercice}[1]{Une sortie moins visible}
\begin{questions}
\item Une lampe éclaire et chauffe son environnement. Nommer deux transferts sortants.
\item L’échauffement signifie-t-il une destruction d’énergie ?
\end{questions}
\end{exercice}
\end{minipage}
\par\noindent\begin{minipage}{\linewidth}
\begin{exercice}[1]{Renouvelable ou fossile}
\begin{questions}
\item Classer vent et charbon.
\item Renouvelable signifie-t-il sans conséquence environnementale ?
\end{questions}
\end{exercice}
\end{minipage}
\clearpage\markright{Je m’entraîne}
\rubrique{Je m’entraîne}
\par\noindent\begin{minipage}{\linewidth}
\begin{exercice}[2]{Un bilan de lampe}
\begin{questions}
\item Une lampe reçoit $200\unite J$ et produit $60\unite J$ de transfert utile sans stockage. Calculer les autres transferts.
\item Vérifier le bilan.
\end{questions}
\end{exercice}
\end{minipage}
\par\noindent\begin{minipage}{\linewidth}
\begin{exercice}[2]{Le moteur}
\begin{questions}
\item Un moteur reçoit $500\unite J$, cède $350\unite J$ mécaniquement et $150\unite J$ thermiquement. Le bilan est-il cohérent sans stockage ?
\item Quel effet utile peut-on choisir ?
\end{questions}
\end{exercice}
\end{minipage}
\par\noindent\begin{minipage}{\linewidth}
\begin{exercice}[2]{Un stockage augmente}
\begin{questions}
\item Un système reçoit $80\unite J$ et en cède $30\unite J$. Quelle variation de stockage ?
\item Peut-on imposer entrée égale sortie ici ?
\end{questions}
\end{exercice}
\end{minipage}
\par\noindent\begin{minipage}{\linewidth}
\begin{exercice}[2]{Une cellule photovoltaïque}
\begin{questions}
\item Décrire l’entrée et une sortie utile d’une cellule photovoltaïque.
\item La lumière transporte-t-elle seulement de l’information ?
\end{questions}
\end{exercice}
\end{minipage}
\par\noindent\begin{minipage}{\linewidth}
\begin{exercice}[2]{Un bilan terrestre}
\begin{questions}
\item Sur une période, la Terre reçoit plus d’énergie qu’elle n’en cède. Que devient son stockage ?
\item Pourquoi un simple dessin ne donne-t-il pas une valeur chiffrée ?
\end{questions}
\end{exercice}
\end{minipage}
\clearpage\markright{Je raisonne}
\rubrique{Je raisonne}
\par\noindent\begin{minipage}{\linewidth}
\begin{exercice}[3]{Comparer deux lampes}
\begin{questions}
\item Pour un même éclairage, A reçoit $100\unite J$ et B $40\unite J$ sur la même durée. Calculer l’économie avec B.
\item Que faut-il vérifier pour que la comparaison soit équitable ?
\end{questions}
\end{exercice}
\end{minipage}
\par\noindent\begin{minipage}{\linewidth}
\begin{exercice}[3]{Une énergie cinétique fournie}
\begin{questions}
\item Pour préparer la 3e, on donne $E_c=\frac12mv^2$. Avec $m=2\unite{kg}$ et $v=3\unite{m/s}$, calculer $E_c$.
\item Si la vitesse double, l’énergie double-t-elle ?
\end{questions}
\end{exercice}
\end{minipage}
\par\noindent\begin{minipage}{\linewidth}
\begin{exercice}[3]{Le mot « perte »}
\begin{questions}
\item Pourquoi les pertes d’un moteur peuvent-elles rester importantes malgré la conservation de l’énergie ?
\item Peut-on récupérer toute énergie aussi facilement ?
\end{questions}
\end{exercice}
\end{minipage}
\par\noindent\begin{minipage}{\linewidth}
\begin{exercice}[3]{Un raccourci sur le climat}
\begin{questions}
\item « Les gaz à effet de serre empêchent toute énergie de sortir. » Corriger.
\item Quelle idée de bilan faut-il retenir ?
\end{questions}
\end{exercice}
\end{minipage}
\par\noindent\begin{minipage}{\linewidth}
\begin{exercice}[3]{Une comparaison incomplète}
\begin{questions}
\item Un appareil émet peu de CO2 pendant son usage. Peut-on en déduire qu’il n’a aucun impact ?
\item Quels postes ajouter ?
\end{questions}
\end{exercice}
\end{minipage}
\pcPage{Défi du chapitre}{L’énergie : transferts et conversions}
\phantomsection\label{defi-ex-12}
\pcDefiTitre{12}{Le convertisseur mystérieux}
\begin{questions}
\item Un appareil reçoit $800\unite J$, cède $500\unite J$ mécaniquement et $220\unite J$ thermiquement. Calculer la variation d’énergie stockée.
\item Un élève affirme qu’il a « perdu 80 J ». Corriger.
\item Comment tester l’hypothèse d’un stockage plutôt qu’un transfert oublié ?
\end{questions}
\pcGrille
